% filepath: PID-in-Modern-Cherenkov/Sections/current-challenges.tex
% Content for the "Current Challenges & Limitations" subsection

Much of modern research is placed into improving photosensor performance in high-rate collider environments. 
Along with the new age SiPMs, improvements are also being made to MCPs, both towards better spatial and temporal resolution. A process known as 
atomic layer deposition (ALD) allows for extremely precise manufacturing of the capillaries in an MCP \nolinebreak\cite{ald-mcps-Popecki_2016}, 
which have been shown to significantly improve sensor ageing \nolinebreak\cite{2013NIMPA.732..388C}. Research into Large Area Picosecond Photodetectors (LAPPDs)
are attempting to reduce the cost per area of traditional MCP-PMTs while matching their performance characteristics. LAPPDs themselves have been shown to have 
good spatial and temporal resolution \nolinebreak\cite{LYASHENKO2020162834}, and further characterisations of HRPPDs, a high-rate upgrade for LAPPDs, 
are ongoing \nolinebreak\cite{LYASHENKO2026170964,lappd-foulds-holt_2025}.

On the software side developments being researched are as described in Sect \nolinebreak\ref{sec:reconstruction-methods} for DIRC systems. 
ML has the potential to offer fast particle classification in all Cherenkov imaging systems, either alone or in conjunction with other PID methods to decrease compute time. 
Effort is being focused into reducing model size, whilst retaining PID performance. This makes the models easier to deploy on the beamline 
and improves classification speed. 

On the electronics side, developments are increasingly driven by the demand for precise timing and high-rate readout to fully exploit modern
photosensors. Fast sampling ASICs\cite{Le-Pottier_2025} and high-bandwidth digitisation schemes\cite{WU2023167994} are required to preserve the picosecond scale 
timing resolution of MCP-based detectors. In high-rate environments, advances in front-end electronics and data acquisition systems aim to 
reduce dead time and manage large data volumes while maintaining low noise performance. Close integration of sensors with dedicated readout 
electronics is therefore an important area of research.

Finally, reducing the cost of detector systems will always be of great interest. The ease of manufacturing for SiPMs offer the potential
to significantly reduce the price tag on detector arrays. Reuse of retired detector parts is also a potential cost saving measure; 
the recent transport of old BaBar boxes for use in the hpDIRC at Brookhaven is one example of successful reuse \nolinebreak\cite{hpDIRC-Greg_2025}.


